\section{引言}

近年来，流式四维重建方法已经能够在顺序处理单目视频的过程中，联合恢复人体运动、相机轨迹和场景几何
~\citep{wang2025cut3r,chen2026human3r,sur2026unicon3r}。这些方法将历史观测压缩到递归状态中，从而实现高效的在线推理，但其建模通常假设相邻帧来自一次连续拍摄。经过剪辑的视频经常违背这一假设：它们由时间上依次出现的多个镜头组成，镜头之间的切换可能突然改变视角、消除视觉重叠，并显著改变可见场景与人体外观
~\citep{pavlakos2022multishot,zhang2025humanmm,on2026multithumbs}。因此，在边界之前积累的历史可能不再适合作为新镜头的递归上下文。本文研究这种非连续输入下的流式重建。这里的“多镜头”是指单目视频中依次出现的镜头片段，而不是同一时刻的同步多视角观测。我们的目标是将每个镜头中的相机轨迹、重建人体和场景几何表示在一个统一世界坐标系中；该坐标系由第一镜头初始化，并作为整段序列共享的重建参考系。

现有方法只解决了这一问题的一个方面。多镜头人体重建方法通过联合处理多个已有镜头或进行序列级优化，在镜头切换前后建立人物关系
~\citep{pavlakos2022multishot,zhang2025humanmm,on2026multithumbs}。这些方法能够有效恢复跨镜头结构，但通常并非为仅依靠有限递归状态、且不访问未来帧的逐步预测而设计。有状态在线重建方法提供了与之互补的能力：它们按照时间顺序处理输入，并避免逐视频的全局优化
~\citep{wang2025cut3r,chen2026human3r,chen2025long3r}。然而，其递归状态被设计为在连续观测中逐步演化，并没有显式处理历史有效性在镜头边界处突然改变的情况。因此，这两类方法都没有直接满足本文的核心需求：在保留先前镜头所建立空间参照的同时，避免将不兼容的历史继续传播到新镜头的重建过程中。

这一缺口源于递归历史在镜头边界处的使用冲突。跨越边界继续传播状态，可以保留先前观测所建立的坐标参照，却也迫使新镜头与在非连续视角下积累的历史发生交互；切换处引入的误差因而可能在后续递归更新中持续传播。重新初始化状态可以消除这种干扰，使新镜头内部恢复连贯的递归重建，但同时也会抹去将该镜头定位到统一世界坐标系所需的参照。因此，关键问题并不是应该保留还是丢弃历史，而是历史的哪些作用应当跨越镜头边界继续存在。我们的核心观察是，历史只应作为临时的对齐参照继续发挥作用：它可以被读取以建立镜头间关系，但不应作为新镜头的递归状态继续传播。使用相同重建模型进行的受控比较支持这一判断：状态延续造成的是随时间变化的相机漂移，而不只是一次固定偏移；独立重建则能够保持新镜头内部的连贯性，却失去了与先前世界坐标系的空间关系（第~\ref{sec:ablation-zh} 节）。

为解决这一冲突，我们提出 \method{}：一个为递归历史的两种作用赋予不同信息生命周期的流式重建框架。在每个镜头边界处，\method{} 临时读取先前状态以推断镜头间的空间关系，但不会将这一过程产生的状态保留为后续重建所使用的记忆。新镜头内部的递归过程从独立初始化的状态开始，并且只有该状态会继续向后传播。在建立空间关系之后，基于几何的人物关联维持镜头边界两侧的身份连续性，相机—人体共享变换则将新镜头的重建结果表示到统一世界坐标系中。\method{} 不需要访问未来帧、修改已经输出的估计结果或执行逐视频优化。在三个多人体基准上，相比连续传播历史状态，\method{} 将按数据集等权计算的平均 W-MPJPE 降低了 9.5\%，并一致改善了相机精度和身份连续性；在 EgoBody 预定义的极端视角子集上，W-MPJPE 的降幅达到 29.9\%。

本文的主要贡献如下：

\begin{itemize}
  \item 我们识别并刻画了递归历史在镜头边界处的使用冲突：历史是建立相邻镜头关系所必需的参照，但将其继续传播到新镜头可能引入持续的重建误差，而将其完全丢弃又会失去镜头间的空间参照。
  \item 我们提出 \method{}，一个将“读取历史以进行对齐”与“传播状态以完成重建”分开处理的流式多镜头重建框架。\method{} 仅将边界之前的历史保留为临时参照，并通过独立初始化的递归状态推进新镜头的重建，同时维持镜头边界两侧的空间与身份一致性。
  \item 通过使用相同重建模型的受控比较以及三个多人体基准上的评测，我们证明这种分离能够改善世界坐标人体重建、相机精度和身份连续性，并在大幅视角变化下取得最明显的收益，同时不需要访问未来帧或执行逐视频优化。
\end{itemize}
